\documentclass[11pt,a4paper]{article}
\usepackage[T1]{fontenc}
\usepackage[utf8]{inputenc}
\usepackage[margin=26mm,headheight=14pt]{geometry}
\usepackage{amsmath,amsthm,mathtools}
\usepackage{newtxtext,newtxmath}
\usepackage{microtype,booktabs,array,longtable}
\usepackage[dvipsnames]{xcolor}
\usepackage{enumitem,fancyhdr,xurl}
\usepackage[colorlinks=true,linkcolor=MidnightBlue,citecolor=MidnightBlue,
 urlcolor=MidnightBlue,pdftitle={Approximate Phase Rigidity and Arithmetic Quadrature Complexity},
 pdfauthor={Research manuscript prepared with ChatGPT}]{hyperref}
\usepackage[capitalise,noabbrev]{cleveref}
\numberwithin{equation}{section}
\setlength{\parskip}{3pt}
\setlength{\parindent}{1.2em}
\setlist{itemsep=3pt,topsep=5pt}
\pagestyle{fancy}\fancyhf{}
\fancyhead[L]{\small Approximate phase rigidity}
\fancyhead[R]{\small Fabius--Rvachev quadrature}
\fancyfoot[C]{\thepage}
\renewcommand{\headrulewidth}{0.3pt}
\theoremstyle{plain}
\newtheorem{theorem}{Theorem}[section]
\newtheorem{lemma}[theorem]{Lemma}
\newtheorem{proposition}[theorem]{Proposition}
\newtheorem{corollary}[theorem]{Corollary}
\theoremstyle{definition}
\newtheorem{definition}[theorem]{Definition}
\newtheorem{example}[theorem]{Example}
\newtheorem{question}[theorem]{Question}
\theoremstyle{remark}
\newtheorem{remark}[theorem]{Remark}
% ed. (2026-09-30): unnumbered environment for editorial notes.
\newtheorem*{ednote}{Editorial note (ProveIt, 2026-09-30)}
\newcommand{\R}{\mathbb R}\newcommand{\C}{\mathbb C}
\newcommand{\N}{\mathbb N}\newcommand{\Z}{\mathbb Z}
\newcommand{\T}{\mathbb R/\mathbb Z}
\newcommand{\ii}{\mathrm i}\newcommand{\dd}{\,\mathrm d}
\newcommand{\sinc}{\operatorname{sinc}}
\newcommand{\supp}{\operatorname{supp}}
\newcommand{\ord}{\operatorname{ord}}
\newcommand{\dist}{\operatorname{dist}}
\newcommand{\vTwo}{\nu_2}
\newcommand{\norm}[1]{\left\lVert#1\right\rVert}
\newcommand{\abs}[1]{\left|#1\right|}
\newcommand{\Err}{\mathcal E}
\newcommand{\HH}{\mathsf H}
\newcommand{\one}{\mathbf 1}
\newcommand{\snapshot}{\texttt{cb1646dc442724f8e298cf259195b6c308367d80}}

\begin{document}
\hypersetup{pageanchor=false}
\begin{titlepage}
\centering
\vspace*{8mm}
{\small\scshape Analysis / Approximation / Arithmetic\par}
\vspace{16mm}
{\LARGE\bfseries Approximate Phase Rigidity\par}
\vspace{3mm}
{\LARGE\bfseries and Arithmetic Quadrature Complexity\par}
\vspace{10mm}
{\large Sharp rational-resonance exponents,\par}
\vspace{2mm}
{\large log-square phase costs, and non-H\"older recovery\par}
\vspace{12mm}
{Research manuscript prepared with ChatGPT\par}
{For Vladimir Reshetnikov\par}
{30 September 2026\par}
\vspace{10mm}
\begin{minipage}{0.94\textwidth}
\begin{abstract}
Exact phase averaging for the Rvachev up-function has an arithmetic rigidity:
at a rational mesh it forces a finite cyclic symmetry, whereas at an irrational
mesh it forces Haar measure. We determine what survives when exactness is
replaced by a small uniform error.

Our principal result gives the exact order of the optimal error near every
positive rational mesh, for every fixed phase budget and every fixed polynomial
degree. Write $M_0=a/b$ in lowest terms, $d=\nu_2(a)$, and let $N\ge b$.
For positive, real signed, and complex mass-one filters with at most $N$ atoms,
the optimal error on the first $r+1$ monomials has order
$|M-M_0|^{\max\{0,d+1+\lfloor\log_2(N/b)\rfloor-r\}}$.
For $N<b$ the error stays bounded away from zero. These lower bounds allow the
nodes and unbounded signed weights to depend on $M$. A regular polygon attains
the order, and its leading response is an explicit rescaled coarse-mesh profile.

For an irrational mesh satisfying a finite-type Diophantine lower bound, the
optimal positive $N$-phase error is $\exp[-\Theta((\log N)^2)]$. Equivalently,
accuracy $\varepsilon$ requires $\exp[\Theta(\sqrt{\log(1/\varepsilon)})]$
phases. We also obtain a finite-frequency Wasserstein inversion inequality.
It yields a matching stretched-logarithmic scale of approximate Haar rigidity,
while excluding every local H\"older inverse estimate. Finally, specially
constructed Liouville meshes permit arbitrarily small errors along prescribed
subsequences and arbitrarily slow inverse moduli. All results have ordinary
proofs; the package includes exact arithmetic checks and numerical diagnostics.
\end{abstract}
\end{minipage}
\vfill
\begin{minipage}{0.94\textwidth}\small
\textbf{Status and scope.} Unrefereed research manuscript. The local optimal
exponent answers the fixed-budget near-rational question in the inspected
ProveIt phase-averaging report. The log-square statements identify the scale,
not an optimal leading constant. No new Lean or Rocq verification is claimed,
and worldwide novelty has not been certified.
\end{minipage}
\end{titlepage}
\hypersetup{pageanchor=true}
\begingroup\small\setlength{\parskip}{0pt}\tableofcontents\endgroup
\newpage

\section{The problem, the repository baseline, and the new results}
\label{sec:intro}

\subsection{Why exact rigidity is not a stability theorem}

A phase-averaged lattice quadrature rule can be exact for every overall
translation only if its phase measure annihilates the appropriate Fourier
aliases. For the Rvachev up-function those aliases have a particularly explicit
zero set. The ProveIt reports \emph{Rigidity and Quantitative Limits of Phase
Averaging} and \emph{Complete Superconvergence Phases and Optimal Phase Filters}
classify exact filters and obtain finite-support obstructions
\cite{repoRigidity,repoFilters}. At an irrational mesh, even uniform exactness
on constants forces continuous Haar averaging. At a rational mesh $a/b$, the
symmetry required through polynomial degree $r$ has order
\begin{equation}\label{eq:oldthreshold}
 L_r=b\,2^{\max(0,r-\nu_2(a))}.
\end{equation}
These are inherited results, not discoveries of this article.

The same rigidity report asks how the quantitative obstruction degenerates
as an irrational mesh approaches a rational scale, and how positive and signed
filters compare away from exactness. These questions are not answered by a
classification at a single parameter value. In particular, an optimizing
filter may move its nodes, change its weights, and, when signs are allowed,
let those weights become arbitrarily large as the mesh changes.

There is a second gap between exactness and approximation. If Haar measure is
the only exact phase measure, does an almost-exact rule have to be close to
Haar measure at an effective rate? The answer is yes for each fixed irrational
mesh, but never at a H\"older rate. Arithmetic assumptions determine whether
one can give a useful uniform scale.

\subsection{The three main conclusions}

The local result, \cref{thm:resonance}, gives a complete answer at every fixed
budget $N$. The exponent depends only on the reduced denominator of the nearby
rational mesh, the two-adic valuation of its numerator, the budget, and the
requested degree. Its proof combines a finite Fourier-moment obstruction with
orders of vanishing of a finite list of multipliers. Thus it remains valid
without any compactness hypothesis on signed weights.

The global irrational result, \cref{thm:complexity}, is different in nature:
$N$ tends to infinity at a fixed irrational mesh. Positivity upgrades an
exponentially weak moment obstruction to a sharp $1/N$ obstruction. Together
with explicit two-sided Fourier estimates, it proves
\begin{equation}\label{eq:scale-intro}
 -\log\Err^+_{N,r}(M)=\Theta((\log N)^2)
\end{equation}
under a finite-type Diophantine hypothesis. The constants in the $\Theta$
notation depend on the fixed mesh, degree, and arithmetic parameters.

The inverse result, \cref{thm:inversion}, bounds the Wasserstein distance of a
phase probability measure from Haar measure using finitely many nonzero
multipliers. This gives a precise stability mechanism and a precise limitation:
\begin{equation}\label{eq:modulus-intro}
 \omega_{M,r}(\varepsilon)
 =\exp[-\Theta(\sqrt{\log(1/\varepsilon)})]
\end{equation}
for finite-type irrational meshes. There is no uniform analogue over arbitrary
irrational meshes. \Cref{thm:liouville} constructs examples with arbitrarily
slow recovery, despite exact uniqueness at zero error.

\subsection{Provenance and the meaning of novelty}

The repository inspection was anchored to commit
\begin{center}\small\snapshot.\end{center}
The two phase-filter reports above were read together with the introduction
of the subsequent fixed-phase maximality report \cite{repoFixed}. The latter
concerns a different problem: a single selected phase at an integer mesh.
We do not repackage its dyadic maximality theorem here.

The Fourier product, compactly supported smooth density, and zero
multiplicities are classical; Arias de Reyna gives the relevant constructions
and formulas \cite{arias}. The annihilating-polynomial obstruction is already
in the repository's rigidity report. The positive-moment argument below is
an elementary Fej\'er-kernel argument, and the inversion mechanism belongs to
the general family of Fourier smoothing arguments. Wasserstein inversion in
deconvolution also appears in a different, statistical setting in the work of
Rousseau and Scricciolo \cite{rousseau}. No claim is made that deconvolution,
Fej\'er kernels, or severe inverse instability are new ideas.

The contributions proposed here are the complete local optimal exponent,
its explicit leading-profile realization, the positive-filter phase-cost
scale for this kernel under a stated arithmetic assumption, and their
combination with a quantified inverse theorem and a Liouville obstruction.
The supplied proofs are independent of unverified assertions elsewhere in the
repository. The source review was targeted, not an audit of the entire corpus
or an exhaustive priority search.

\section{Kernel, filters, and the error being optimized}
\label{sec:setup}

\subsection{Normalization}

Let $V_1,V_2,\ldots$ be independent uniform random variables on $[-1,1]$ and
set $Y=\sum_{j\ge1}2^{-j}V_j$. Its law has the normalized Rvachev density
$u\in C_c^\infty(\R)$, supported on $[-1,1]$, with $u\ge0$, $u$ even, and
$\int u=1$. In the convention
\begin{equation}\label{eq:fourier}
 \widehat f(t)=\int_\R f(x)e^{-2\pi\ii tx}\dd x,
\end{equation}
its Fourier transform is
\begin{equation}\label{eq:product}
 H(t)=\widehat u(t)=\prod_{j=0}^{\infty}
 \frac{\sin(\pi t/2^j)}{\pi t/2^j},
\end{equation}
where a factor at zero is interpreted as one. This agrees with the
normalization in \cite{arias} and the inspected phase reports.

The product converges locally uniformly in the complex plane: after finitely
many factors, the difference of a factor from one is bounded by a constant
multiple of $4^{-j}$ on each compact set. It defines an entire function.
Since $u\in C_c^\infty$, every derivative of $H$ is rapidly decreasing on the
real axis. These facts also follow directly from the product; a quantitative
version is proved in \cref{sec:envelopes}.

\begin{lemma}[Zero order and leading derivative]\label[lemma]{lem:zeros}
The real zeros of $H$ are exactly the nonzero integers. If
$t_0=2^v m\ne0$ with $m$ odd and $p=v+1$, then
\begin{equation}\label{eq:leadingjet}
 \ord_{t_0}H=p,
 \qquad H^{(p)}(t_0)=-\frac{p!}{t_0^p}H(m/2)\ne0.
\end{equation}
\end{lemma}
\begin{proof}
Exactly the factors indexed by $0,\ldots,v$ vanish at $t_0$. Their linear
coefficients in $t-t_0$ are $(-1)^{t_0/2^j}/t_0$. The product of these signs is
$(-1)^{m(2^{v+1}-1)}=-1$. The remaining product at $t_0$ is $H(m/2)$, which is
nonzero: no factor there vanishes, and the tail is an absolutely convergent
product of factors tending to one. This proves the formula. The same
nonvanishing-product argument excludes any other real zero.
\end{proof}

% ed. (2026-09-30): uncited Lean declarations proving lem:zeros.
\begin{ednote}
\Cref{lem:zeros} is machine-checked in the repository, for the same product
\eqref{eq:product} (\path{Fabius.rvachevFourierProduct}) and even over
$\C$, although this article does not cite it: the complex zeros of $H$ are
exactly the nonzero integers
(\path{Fabius.rvachevFourierProduct_eq_zero_iff},
\path{Analysis/FabiusFunction/Lean/FabiusFunction/FourierProduct.lean});
at a nonzero integer $t_0$ the analytic order is $\nu_2(|t_0|)+1$
(\path{Fabius.analyticOrderAt_rvachevFourierProduct_int}), and the
derivative of that order is $-p!\,H(m/2)/t_0^{p}$ with $m$ the odd part of
$|t_0|$, which is nonzero
(\path{Fabius.iteratedDeriv_rvachevFourierProduct_int} and
\path{Fabius.iteratedDeriv_rvachevFourierProduct_int_ne_zero}), all in
\path{Analysis/FabiusFunction/Lean/FabiusFunction/IntegerZeroAnalyticOrder.lean}.
The zero-multiplicity and leading-derivative half of the ``Kernel jets''
layer in \cref{sec:formal} is therefore already formalized. No theorem of
this article has a Lean statement.
\end{ednote}

\subsection{The translation-uniform error}

Write $\lambda$ for normalized Haar measure on $\T=\R/\Z$, and
$m_j=\int x^ju(x)\dd x$. For $M>0$ define
\begin{equation}\label{eq:Q}
 Q_{M,\theta}(P)=\frac1M\sum_{k\in\Z}
 P\!\left(\frac{k+\theta}{M}\right)
 u\!\left(\frac{k+\theta}{M}\right).
\end{equation}
The physical spacing is $1/M$. This is a finite sum, because $u$ is compactly
supported. A phase filter is a finite complex Borel measure $\nu$ on $\T$ of
mass one. We use the convolution convention
\begin{align}
 Q_{M,\nu,\theta}(P)&=\int_\T Q_{M,\theta-\phi}(P)\dd\nu(\phi),
 \label{eq:averaged}\\
 D_j(M,\nu;\theta)&=Q_{M,\nu,\theta}(x^j)-m_j,
 \qquad
 E_{M,r}(\nu)=\max_{0\le j\le r}\norm{D_j(M,\nu;\cdot)}_\infty.
 \label{eq:error}
\end{align}
Using $\theta+\phi$ instead merely reflects the phase measure and does not
change any optimized quantity. The norm in \eqref{eq:error} is also the
operator norm for polynomials $\sum_{j=0}^r c_jx^j$ with
$\sum_j|c_j|\le1$, taking the supremum over $\theta$.

Let $\mathcal A_N^+$, $\mathcal A_N^{\R}$, and $\mathcal A_N^{\C}$ denote,
respectively, the positive, real signed, and complex mass-one filters with
at most $N$ distinct nonzero atoms. Define
\begin{equation}\label{eq:optimal}
 \Err_{N,r}^{\mathbb K}(M)=
 \inf_{\nu\in\mathcal A_N^{\mathbb K}} E_{M,r}(\nu),
 \qquad \mathbb K\in\{+,\R,\C\}.
\end{equation}
Duplicate phases are coalesced; zero weights do not count. In particular,
\begin{equation}\label{eq:classes}
 \Err_{N,r}^{\C}\le\Err_{N,r}^{\R}\le\Err_{N,r}^{+}.
\end{equation}
There is no bound on total variation in either signed class.

\begin{lemma}[Alias identity and smooth parameter dependence]\label[lemma]{lem:alias}
Put $\widehat\nu(n)=\int e^{-2\pi\ii n\phi}\dd\nu(\phi)$. Then
\begin{equation}\label{eq:alias}
 D_j(M,\nu;\theta)=
 \sum_{n\ne0}\left(\frac{\ii}{2\pi}\right)^j
 H^{(j)}(nM)\widehat\nu(n)e^{2\pi\ii n\theta}.
\end{equation}
For a fixed finite measure, the series and every fixed number of its
$M$- and $\theta$-derivatives converge uniformly when $M$ ranges over a compact
subinterval of $(0,\infty)$.
\end{lemma}
\begin{proof}
The Fourier transform of $x^ju(x)$ is $(\ii/(2\pi))^jH^{(j)}$.
Poisson summation for this smooth compactly supported function gives the
unfiltered identity. Its zero frequency is $m_j$. Integrating against
$\nu$ gives \eqref{eq:alias}. A derivative contributes only a fixed power of
$n$ and a higher derivative of $H$. Rapid decrease, together with
$|\widehat\nu(n)|\le\norm{\nu}_{\mathrm{TV}}$, proves uniform convergence.
\end{proof}

\subsection{Regular polygons and exact refinement}

Let
\begin{equation}\label{eq:grid}
 \nu_L=\frac1L\sum_{h=0}^{L-1}\delta_{h/L}.
\end{equation}
Its Fourier coefficients are one on $L\Z$ and zero elsewhere. Reindexing
\eqref{eq:Q}, or using \eqref{eq:alias}, gives
\begin{equation}\label{eq:refinement}
 Q_{M,\nu_L,\theta}(P)=Q_{LM,L\theta}(P).
\end{equation}
Thus $L$ equally weighted phase shifts are exactly an $L$-fold mesh refinement,
not a new choice of arbitrary physical weights.

For completeness, \cref{lem:zeros,lem:alias} recover the inherited exact
classification. At $M=a/b$ in lowest terms, all multipliers through degree
$r$ vanish at frequency $n\ne0$ precisely when $L_r\mid n$, with $L_r$ as in
\eqref{eq:oldthreshold}. Therefore exactness means that $\widehat\nu$ is
supported on $L_r\Z$, equivalently that $\nu$ is invariant under rotation by
$1/L_r$. Fourier uniqueness follows by approximating continuous test functions
by trigonometric polynomials. At an irrational $M$, $H(nM)\ne0$ for every
$n\ne0$, so exactness even for $j=0$ forces $\nu=\lambda$.

% ed. (2026-09-30): relation to the canonical comb synthesis and its Lean
% counterpart, which the article does not cite.
\begin{ednote}
For an integer mesh $M_0=a$ and the one-atom filter $\nu=\delta_0$, the
classification above says that the unfiltered comb is exact through degree
$r$ at every phase if and only if $r\le\nu_2(a)$. This is Theorem
\texttt{thm:up-composite-mesh} of the canonical synthesis
\path{inverse-and-sampling/comb-interpolation/comb_interpolation_synthesis/}
(part \emph{Additive dyadic foundations}, same normalization
\eqref{eq:Q}) together with its sharpness, machine-checked in the unscaled
form (nodes $k+\theta$, kernel $u(x/M)$) as
\path{Fabius.rvachevCombExactThrough_iff_padicValNat}
(\path{Analysis/FabiusFunction/Lean/FabiusFunction/CompositeMeshSharpness.lean}).
The phase results of that part (\texttt{thm:phase-average},
\texttt{thm:phase-cubature}, \texttt{thm:tm-bernoulli-egf},
\texttt{thm:root-unity-filter}, \texttt{thm:phase-zero-set}) are exact
identities at integer meshes; the synthesis contains no statement at
irrational meshes and no approximate, Haar-measure or Wasserstein statement.
\Cref{thm:resonance,thm:profile,thm:complexity,thm:inversion,thm:liouville}
therefore have no counterpart there.
\end{ednote}

\section{Finite Fourier witnesses: signs versus positivity}
\label{sec:witnesses}

Define the finite multiplier barrier
\begin{equation}\label{eq:barrier}
 B_{n,r}(M)=\max_{0\le j\le r}
 \frac{|H^{(j)}(nM)|}{(2\pi)^j},
 \qquad b_{N,r}(M)=\min_{1\le n\le N}B_{n,r}(M).
\end{equation}
The following elementary bounds will be used in two very different limits:
$M\to a/b$ with $N$ fixed, and $N\to\infty$ with $M$ fixed.

\begin{lemma}[A mass-one atomic measure has a low Fourier witness]
\label[lemma]{lem:momentbarrier}
For every $\nu\in\mathcal A_N^{\C}$,
\begin{equation}\label{eq:complexmoment}
 \max_{1\le n\le N}|\widehat\nu(n)|\ge\frac1{2^N-1}.
\end{equation}
For every $\nu\in\mathcal A_N^{+}$, the stronger sharp inequality holds:
\begin{equation}\label{eq:positivemoment}
 \max_{1\le n\le N}|\widehat\nu(n)|\ge\frac1N.
\end{equation}
\end{lemma}
\begin{proof}
Suppose first that the actual support has $K\le N$ points. Put
$z_h=e^{-2\pi\ii\phi_h}$ and
$A(z)=\prod_{h=1}^K(z-z_h)=\sum_{j=0}^K a_jz^j$.
Since $A$ vanishes on the support,
$\sum_{j=0}^K a_j\widehat\nu(j)=0$. Here $|a_0|=1$ and
$\widehat\nu(0)=1$. The elementary-symmetric coefficient bound gives
$\sum_{j=1}^K|a_j|\le2^K-1$. This proves \eqref{eq:complexmoment} with $K$,
and hence with $N$. It uses neither positivity nor a bound on the weights.
This is the obstruction already used in \cite{repoRigidity}.

For positive weights $w_h$, use the Fej\'er kernel
\begin{equation}\label{eq:fejer}
 F_N(x)=\sum_{|n|\le N}\left(1-\frac{|n|}{N+1}\right)e^{2\pi\ii nx}
 =\frac1{N+1}\left|\sum_{k=0}^Ne^{2\pi\ii kx}\right|^2.
\end{equation}
It is pointwise nonnegative and $F_N(0)=N+1$. Therefore
\begin{align*}
 1+2\sum_{n=1}^N\left(1-\frac n{N+1}\right)|\widehat\nu(n)|^2
 &=\sum_{h,k}w_hw_k F_N(\phi_h-\phi_k)\\
 &\ge(N+1)\sum_hw_h^2\ge\frac{N+1}{N}.
\end{align*}
Since $2\sum_{n=1}^N(1-n/(N+1))=N$, the largest squared Fourier coefficient
is at least $1/N^2$.

Sharpness of \eqref{eq:positivemoment} is witnessed by the equally weighted
measure on any $N$ vertices of a regular $(N+1)$-gon. For $1\le n\le N$, its
Fourier coefficient is minus the omitted vertex's $n$th power divided by $N$.
\end{proof}

\begin{corollary}[Explicit error certificates]\label[corollary]{cor:barrier}
For every $M>0$, $r\ge0$, and $N\ge1$,
\begin{equation}\label{eq:errorbarriers}
 \Err_{N,r}^{\C}(M)\ge\frac{b_{N,r}(M)}{2^N-1},
 \qquad
 \Err_{N,r}^{+}(M)\ge\frac{b_{N,r}(M)}N.
\end{equation}
\end{corollary}
\begin{proof}
The absolute value of any Fourier coefficient of $D_j$ is at most its
uniform norm. By \eqref{eq:alias},
$E_{M,r}(\nu)\ge B_{n,r}(M)|\widehat\nu(n)|$ for each $n$.
Take the frequency supplied by \cref{lem:momentbarrier}.
\end{proof}

\begin{remark}
The $1/N$ estimate is sharp as a statement about low moments. It does not
assert that the resulting quadrature lower bound is sharp in its leading
constant. Also, positivity cannot simply be deleted from its proof: the
nonnegative off-diagonal Fej\'er terms are multiplied by products of weights.
\end{remark}

\section{The exact optimal exponent near every rational mesh}
\label{sec:resonance}

\subsection{Statement of the local theorem}

Fix $M_0=a/b>0$ in lowest terms. Define
\begin{equation}\label{eq:zN}
 z_N(a/b)=
 \begin{cases}
 0,&N<b,\\
 \nu_2(a)+1+\lfloor\log_2(N/b)\rfloor,&N\ge b,
 \end{cases}
 \qquad q_{N,r}(a/b)=\max\{0,z_N(a/b)-r\}.
\end{equation}
The first case is essential: a large two-adic numerator does not remove the
need to pay for the denominator.

\begin{theorem}[Sharp rational-resonance exponent]\label{thm:resonance}
Let $a,b,N$ be positive integers with $\gcd(a,b)=1$, and let $r\ge0$.
There are constants $c,C,\eta>0$, depending on these fixed data, such that
for $0<|M-a/b|<\eta$,
\begin{equation}\label{eq:resonance}
 c|M-a/b|^{q_{N,r}(a/b)}
 \le\Err_{N,r}^{\C}(M)
 \le\Err_{N,r}^{\R}(M)
 \le\Err_{N,r}^{+}(M)
 \le C|M-a/b|^{q_{N,r}(a/b)}.
\end{equation}
When the exponent is zero, this means a positive lower bound, not convergence
to zero. When it is positive, the explicit positive filter
\begin{equation}\label{eq:optimalordergrid}
 L=b\,2^s,\qquad s=\lfloor\log_2(N/b)\rfloor,\qquad \nu_L
\end{equation}
attains the order. The lower bound permits arbitrary mesh-dependent nodes and
arbitrarily large real or complex weights.
\end{theorem}

This is a local statement about the optimal \emph{order}, with $N$ and $r$
fixed. It does not state that the same regular polygon always minimizes the
leading constant, or that the neighborhood and constants are uniform as the
budget grows.

\subsection{The finite list that determines the exponent}

For $1\le n\le N$ let $v_n$ be the zero order of $H$ at $nM_0$, with order
zero assigned to a nonzero value. By \cref{lem:zeros},
\begin{equation}\label{eq:vn}
 v_n=
 \begin{cases}
 0,&b\nmid n,\\
 d+1+\nu_2(n/b),&b\mid n,
 \end{cases}
 \qquad d=\nu_2(a).
\end{equation}
Consequently
\begin{equation}\label{eq:maxvn}
 \max_{1\le n\le N}v_n=z_N(a/b).
\end{equation}
Indeed, for $N\ge b$, the largest two-adic valuation of an integer not exceeding
$N/b$ is $\lfloor\log_2(N/b)\rfloor$, attained at a power of two.

\begin{lemma}[Local multiplier asymptotics]\label[lemma]{lem:localbarrier}
If $v_n>r$, then, as $\delta\to0$,
\begin{equation}\label{eq:localB}
 B_{n,r}(M_0+\delta)=
 \frac{|H^{(v_n)}(nM_0)|n^{v_n-r}}
 {(2\pi)^r(v_n-r)!}\,|\delta|^{v_n-r}(1+o(1)).
\end{equation}
If $v_n\le r$, then $B_{n,r}(M_0)>0$ and $B_{n,r}$ is continuous there.
In particular,
\begin{equation}\label{eq:b-asymptotic}
 b_{N,r}(M_0+\delta)\asymp |\delta|^{q_{N,r}(M_0)}.
\end{equation}
\end{lemma}
\begin{proof}
Taylor expansion at $nM_0$ shows that $H^{(j)}(nM_0+n\delta)$ has order
$v_n-j$ when $j<v_n$. If $v_n>r$, the $j=r$ term in the maximum has the lowest
order, with the nonzero coefficient in \eqref{eq:localB}; all terms with
$j<r$ are smaller by a positive power of $|\delta|$. If $v_n\le r$, the
$j=v_n$ term is nonzero at the center. Taking the minimum over a finite list
selects its largest vanishing order, giving \eqref{eq:b-asymptotic}.
\end{proof}

For a positive exponent $q$, the leading constant of $b_{N,r}$ is explicitly
\begin{equation}\label{eq:barrierconstant}
 \min_{\substack{1\le n\le N\\v_n=z_N}}
 \frac{|H^{(z_N)}(nM_0)|n^q}{(2\pi)^r q!}.
\end{equation}
Thus the lower bound is more than an abstract compactness argument: it is
controlled by finitely many nonzero half-integer product values through
\eqref{eq:leadingjet}.

\subsection{Proof of the theorem}

\begin{proof}[Proof of \cref{thm:resonance}]
The lower bound follows immediately from
\cref{cor:barrier,lem:localbarrier}; the factor $2^N-1$ is constant in this
local limit. The nesting of the three classes gives the middle inequalities.

Suppose first that $q=q_{N,r}(M_0)>0$. Then $N\ge b$. Let $s,L$ be as in
\eqref{eq:optimalordergrid} and set $p=d+s+1$, so $q=p-r$.
The only nonzero Fourier coefficients of $\nu_L$ occur at frequencies $Ln$.
At the central mesh,
\[
 LnM_0=a2^s n,
\]
which is a nonzero integer with zero multiplicity at least $p$. For
$0\le j\le r$ and $0\le\ell<q$, one has $j+\ell<p$, hence
\[
 \frac{\partial^\ell}{\partial M^\ell}H^{(j)}(LnM)\bigg|_{M=M_0}
 =(Ln)^\ell H^{(j+\ell)}(LnM_0)=0.
\]
By \cref{lem:alias}, these are also the derivatives of the uniformly convergent
alias series. Taylor's theorem in the uniform norm therefore gives
$E_{M_0+\delta,r}(\nu_L)=O(|\delta|^q)$.
This proves the desired upper bound with a positive filter of $L\le N$ atoms.

If $q=0$, choose any one-atom filter. Its error is locally bounded as a
continuous function of $M$, by \cref{lem:alias}. This proves the upper bound
with exponent zero. No minimizing signed filter, or convergent subsequence of
signed filters, has been used anywhere in the proof.
\end{proof}

\begin{corollary}[One doubling, one extra robustness order]\label[corollary]{cor:doubling}
Once $N\ge b$ and the requested degree is exactly reproducible at $M_0$,
the local order increases by one precisely when $N$ crosses a threshold
$b2^s$. Increasing $r$ by one decreases that order by one, until the order
reaches zero. In particular, allowing signs cannot improve the exponent.
\end{corollary}

\begin{example}[An odd numerator]\label{ex:three-fifths}
At $M_0=3/5$, constant reproduction has the following optimal local orders:
\begin{center}
\begin{tabular}{ccl}
\toprule
Phase budget $N$ & Exponent & Consequence\\
\midrule
$1$--$4$ & $0$ & A nonzero error remains at the rational mesh.\\
$5$--$9$ & $1$ & Linear robustness.\\
$10$--$19$ & $2$ & Quadratic robustness.\\
$20$--$39$ & $3$ & Cubic robustness.\\
$40$--$79$ & $4$ & Fourth-order robustness.\\
\bottomrule
\end{tabular}
\end{center}
For degree $r=2$, the first positive local order occurs at $N=20$; the
orders in the last two rows become one and two. The assertions apply to the
best adaptive complex filter as well as to the best positive filter.
\end{example}

\begin{example}[A numerator with extra two-adic divisibility]
At $M_0=4/3$, three through five phases give constant error of exact optimal
order three, while six through eleven phases give order four. With degree
$r=2$, those orders become one and two. Fewer than three phases cannot make
even the constant error vanish near the central mesh.
\end{example}

\section{The leading response is a coarse-mesh profile}
\label{sec:profile}

The upper-bound proof hides a useful explicit formula. Retain $M_0=a/b$,
$a=2^d m$ with $m$ odd, $L=b2^s$, $p=d+s+1$, and assume $r<p$.
Set $q=p-r$ and define the smooth periodic function
\begin{equation}\label{eq:Psi}
 \Psi_{m,r}(x)=
 \sum_{\substack{n\ne0\\n\ \mathrm{odd}}}
 \left(\frac{\ii}{2\pi n}\right)^r H(mn/2)e^{2\pi\ii nx}.
\end{equation}
Rapid decrease makes this series and all its derivatives uniformly
convergent. It is real valued, and its first Fourier coefficient is nonzero.

\begin{theorem}[Explicit leading profile]\label{thm:profile}
Uniformly in $\theta$, as $\delta\to0$,
\begin{align}
 D_r(M_0+\delta,\nu_L;\theta)
 &=-\frac{p!}{q!L^rM_0^p}\,
 \delta^q\Psi_{m,r}(L\theta)+O(|\delta|^{q+1}),
 \label{eq:profile}\\
 D_j(M_0+\delta,\nu_L;\theta)&=O(|\delta|^{p-j})\quad(0\le j<r).
 \label{eq:lowerprofiles}
\end{align}
Consequently
\begin{equation}\label{eq:profile-norm}
 E_{M_0+\delta,r}(\nu_L)=
 \frac{p!\norm{\Psi_{m,r}}_\infty}{q!L^rM_0^p}|\delta|^q
 +O(|\delta|^{q+1}),
\end{equation}
with a strictly positive leading constant.
\end{theorem}
\begin{proof}
In the alias expansion for $D_r$, the coefficient of $\delta^q$ at frequency
$Ln$ is
\[
 \left(\frac{\ii}{2\pi}\right)^r\frac{(Ln)^q}{q!}
 H^{(p)}(a2^s n).
\]
It is zero when $n$ is even, because the multiplicity then exceeds $p$.
For odd $n$, \cref{lem:zeros} makes it
\[
 -\frac{p!}{q!L^rM_0^p}
 \left(\frac{\ii}{2\pi n}\right)^r H(mn/2).
\]
This proves \eqref{eq:profile}; the uniform remainder follows from the next
uniform $M$-derivative in \cref{lem:alias}. The same argument through degree
$p-j-1$ proves \eqref{eq:lowerprofiles}. The uniform norm is Lipschitz with
respect to uniform perturbations, so taking the maximum over $j$ gives
\eqref{eq:profile-norm}. Nonvanishing of the first coefficient implies
$\norm{\Psi_{m,r}}_\infty>0$.
\end{proof}

For constants, the profile has a particularly direct interpretation:
\begin{equation}\label{eq:coarseprofile}
 \Psi_{m,0}(x)=Q_{m/2,x}(1)-1.
\end{equation}
Indeed the even aliases at the mesh $m/2$ are $H(mk)=0$ for $k\ne0$.
Thus the first nonzero response near a large refined rational mesh is a
rescaled defect of the fixed half-integer mesh determined by the odd part
of the numerator. In particular,
\begin{equation}\label{eq:coarsemass}
 D_0(M_0+\delta,\nu_L;\theta)
 =-\left(\frac\delta{M_0}\right)^p
 \bigl(Q_{m/2,L\theta}(1)-1\bigr)+O(|\delta|^{p+1}).
\end{equation}

When $m=1$, this yields an exact leading norm constant:
\begin{equation}\label{eq:dyadicconstant}
 E_{M_0+\delta,0}(\nu_L)
 =\left|\frac\delta{M_0}\right|^p+O(|\delta|^{p+1}).
\end{equation}
To see this without numerical evaluation, write $Y=V_1/2+Z$ with
$|Z|\le1/2$. Convolution with the uniform density on $[-1/2,1/2]$ gives
$0\le u\le1$ and $u(0)=1$. At mesh $1/2$, at most one nonzero interior
sample occurs, so $0\le Q_{1/2,x}(1)\le2$. Its values at $x=0$ and $x=1/2$
are two and zero, respectively. Hence $\norm{\Psi_{1,0}}_\infty=1$.

\begin{remark}[What is and is not optimized]
Equation \eqref{eq:profile-norm} is an asymptotic equivalence for the specified
regular polygon. \Cref{thm:resonance} optimizes its exponent over all allowed
filters. The leading constant for the fully optimized problem may be smaller;
we do not identify it here.
\end{remark}

\section{Two-sided Fourier estimates with explicit envelopes}
\label{sec:envelopes}

\subsection{A uniform upper envelope for every fixed derivative}

For $t\ge1$ put
\begin{equation}\label{eq:envelope}
 G(t)=t^{-1/2}\exp\!\left(-\frac{(\log t)^2}{2\log2}\right).
\end{equation}

\begin{lemma}[Differentiated log-square envelope]\label[lemma]{lem:upperH}
For every integer $r\ge0$ and real $|t|\ge1$,
\begin{equation}\label{eq:upperH}
 \max_{0\le j\le r}\frac{|H^{(j)}(t)|}{(2\pi)^j}
 \le 2^rG(|t|).
\end{equation}
\end{lemma}
\begin{proof}
It suffices to consider $t\ge1$. Write $f(x)=\sin(\pi x)/(\pi x)$.
For $|x|\ge1$, $|f(x)|\le1/|x|$. For $\ell\ge1$, differentiating
$f(x)=\tfrac12\int_{-1}^1e^{\ii\pi xs}\dd s$ and integrating by parts once
shows
\[
 |f^{(\ell)}(x)|\le\frac{2\pi^{\ell-1}}{|x|}
 \le\frac{(2\pi)^\ell}{|x|}.
\]
Let $J=\lfloor\log_2t\rfloor+1$ and split
$H(t)=\prod_{h=0}^{J-1} f(t/2^h)\,H(t/2^J)$.
Each initial argument is at least one. Also, because $|Y|\le1$,
$|H^{(\ell)}(y)|\le(2\pi)^\ell$ for every real $y$.
Leibniz's rule therefore gives
\[
 |H^{(j)}(t)|\le
 (2\pi)^j\left(\sum_{h=0}^J2^{-h}\right)^j
 \frac{2^{J(J-1)/2}}{t^J}
 \le(4\pi)^j\frac{2^{J(J-1)/2}}{t^J}.
\]
Writing $\log_2t=k+\xi$, $0\le\xi<1$, shows
\[
 \log\left(\frac{2^{J(J-1)/2}}{t^J}\right)
 =-\frac{(\log t)^2}{2\log2}-\frac12\log t
   +\frac{\log2}{2}(\xi^2-\xi)
 \le\log G(t).
\]
Divide by $(2\pi)^j$ and take the maximum.
\end{proof}

\begin{corollary}[An explicit regular-grid error bound]\label[corollary]{cor:gridupper}
For $MN\ge2$,
\begin{equation}\label{eq:gridupper}
 E_{M,r}(\nu_N)\le 6\,2^r G(MN).
\end{equation}
\end{corollary}
\begin{proof}
Only aliases $Nn$ survive. By \cref{lem:upperH}, their absolute sum is at most
$2^{r+1}\sum_{n\ge1}G(MNn)$. Moreover,
\[
 \frac{G(xn)}{G(x)}
 =n^{-1/2-\log_2x}
   \exp\!\left(-\frac{(\log n)^2}{2\log2}\right)
 \le n^{-3/2}\quad(x\ge2).
\]
The elementary integral bound gives $\sum_{n\ge1}n^{-3/2}\le3$.
\end{proof}

\subsection{A lower envelope away from arithmetic resonances}

Write $\|t\|_{\Z}=\dist(t,\Z)$. The lower bound must retain this distance:
a global positive lower envelope for $H$ is impossible because $H$ has zeros.

\begin{lemma}[A nonzero product lower bound]\label[lemma]{lem:lowerH}
If $t\ge1$ is not an integer and $J=\lceil\log_2(2t)\rceil$, then
\begin{equation}\label{eq:lowerH}
 |H(t)|\ge c_*\left(\frac{2\|t\|_{\Z}}{\pi t}\right)^J,
 \qquad c_*=e^{-\pi^2/9}.
\end{equation}
More generally, for $0\le s\le1/2$,
\begin{equation}\label{eq:tailbound}
 e^{-4\pi^2s^2/9}\le H(s)\le1.
\end{equation}
\end{lemma}
\begin{proof}
For every real $x$, $|\sin(\pi x)|\ge2\|x\|_{\Z}$. Since
$2^h\Z\subseteq\Z$,
\[
 \left\|\frac t{2^h}\right\|_{\Z}
 =2^{-h}\dist(t,2^h\Z)\ge2^{-h}\|t\|_{\Z}.
\]
Thus each of the first $J$ sinc factors has magnitude at least
$2\|t\|_{\Z}/(\pi t)$.

For the tail let $x_h=\pi s/2^h$. Taylor's elementary lower bound gives
$\sin x_h/x_h\ge1-x_h^2/6$. Here $x_h^2/6\le\pi^2/24<1/2$, and
$\log(1-y)\ge-2y$ on $[0,1/2]$. Summing gives
\[
 \log H(s)\ge-\frac13\sum_{h\ge0}x_h^2
 =-\frac{4\pi^2s^2}{9}.
\]
The upper bound is immediate. Since $t/2^J\le1/2$, this proves
\eqref{eq:lowerH} as well.
\end{proof}

\begin{definition}[Finite-type mesh]\label[definition]{def:type}
An irrational $M>0$ satisfies a finite-type bound if there are constants
$\kappa>0$ and $\tau\ge1$ such that
\begin{equation}\label{eq:type}
 \|nM\|_{\Z}\ge\kappa n^{-\tau}\qquad(n\ge1).
\end{equation}
The constants are part of the hypothesis. We do not assume a uniform
$\kappa$ over all meshes.
\end{definition}

\begin{proposition}[A finite-type spectral floor]\label[proposition]{prop:floor}
Under \eqref{eq:type}, there exist constants $B,C>0$ such that for every
$N\ge2$,
\begin{equation}\label{eq:floor}
 h_M(N):=\min_{1\le n\le N}|H(nM)|
 \ge\exp\!\left[-\frac{\tau+1}{\log2}(\log N)^2-B\log N-C\right].
\end{equation}
\end{proposition}
\begin{proof}
For $nM\ge1$, use \cref{lem:lowerH}. Set
$a_0=\min\{1,2\kappa/(\pi M)\}>0$ and choose a fixed $C_M\ge1$ such that
$\lceil\log_2(2nM)\rceil\le\log_2n+C_M$ for all such $n$.
Then
\[
 |H(nM)|\ge c_*
 \bigl(a_0 n^{-(\tau+1)}\bigr)^{\log_2n+C_M}.
\]
Taking logarithms gives the claimed quadratic leading coefficient and linear
remainder. The finitely many indices with $nM<1$ have nonzero values; enlarge
$C$ to include their minimum. Finally replace $n$ by the upper bound $N$ in
the resulting decreasing lower envelope.
\end{proof}

\begin{example}[A fully elementary finite-type example]
For $M=\sqrt2$, let $k$ be the nearest integer to $n\sqrt2$.
The nonzero integer $|2n^2-k^2|$ is at least one. Also
$n\sqrt2+k<4n$ for $n\ge1$. Hence
\[
 \|n\sqrt2\|_{\Z}
 =\frac{|2n^2-k^2|}{n\sqrt2+k}\ge\frac1{4n}.
\]
Thus \eqref{eq:type} holds with $\kappa=1/4$ and $\tau=1$, without appealing
to a deeper approximation theorem.
\end{example}

\section{Optimal positive phase complexity at irrational meshes}
\label{sec:complexity}

\begin{theorem}[The log-square phase-cost scale]\label{thm:complexity}
Fix $r\ge0$ and an irrational mesh $M$ satisfying \eqref{eq:type}. Then
\begin{equation}\label{eq:log-square}
 \Err^+_{N,r}(M)=\exp[-\Theta((\log N)^2)]\qquad(N\to\infty).
\end{equation}
More explicitly,
\begin{equation}\label{eq:limconstants}
 \frac1{2\log2}
 \le\liminf_{N\to\infty}\frac{-\log\Err^+_{N,r}(M)}{(\log N)^2}
 \le\limsup_{N\to\infty}\frac{-\log\Err^+_{N,r}(M)}{(\log N)^2}
 \le\frac{\tau+1}{\log2}.
\end{equation}
If $\mathsf N_{M,r}(\varepsilon)$ is the smallest positive phase budget whose
optimal error is at most $\varepsilon$, then
\begin{equation}\label{eq:cost}
 \mathsf N_{M,r}(\varepsilon)
 =\exp[\Theta(\sqrt{\log(1/\varepsilon)})]
 \qquad(\varepsilon\downarrow0).
\end{equation}
\end{theorem}
\begin{proof}
By \cref{cor:barrier}, using just the zeroth multiplier,
\begin{equation}\label{eq:basicirrationalbounds}
 \frac{h_M(N)}N\le\Err^+_{N,r}(M)
 \le E_{M,r}(\nu_N)\le6\,2^rG(MN)
\end{equation}
once $MN\ge2$. The lower estimate \eqref{eq:floor} and the explicit formula
for $G$ prove \eqref{eq:log-square} and \eqref{eq:limconstants}.

The positive infimum is attained: represent a filter by $N$ circle points
and a point of the compact probability simplex, allowing repeated nodes and
zero weights. The error is continuous in these parameters, uniformly in
$\theta$. Thus there is a minimum. The optimized error is nonincreasing in
$N$, is positive for each finite $N$ by \eqref{eq:basicirrationalbounds}, and
tends to zero by the grid upper bound. Inverting the two log-square
inequalities proves \eqref{eq:cost}.
\end{proof}

\subsection{Interpretation and limitations}

The result says that regular polygons have the correct \emph{logarithmic
scale} of complexity among all positive filters, even those with highly
nonuniform weights and nonuniformly spaced nodes. It does not say they have
the smallest error at each $N$, nor that the lower and upper constants in
\eqref{eq:limconstants} coincide. For the elementary mesh $M=\sqrt2$, those
constants are $1/(2\log2)$ and $2/\log2$.

The distinction from signed filters is real. For unrestricted complex weights,
our unconditional lower estimate is only
\begin{equation}\label{eq:signed-global}
 \Err^{\C}_{N,r}(M)\ge\frac{h_M(N)}{2^N-1}.
\end{equation}
This is enough to prove the exact fixed-$N$ rational-resonance exponent, but
not enough to prove \eqref{eq:log-square} when $N$ grows. We do not replace
$2^N-1$ by $N$ for signed weights, and do not claim the positive complexity
law for that larger class.

The theorem also must not be confused with a claim about arbitrary Gaussian
quadrature. The rule family retains its kernel-dependent physical weights,
and one filter must work for every overall translation. Removing that
translation quantifier changes the optimization problem fundamentally.

\section{Approximate Haar rigidity and a finite-frequency inverse bound}
\label{sec:inverse}

For probability measures on the circle, define the Wasserstein distance by
its usual dual formula
\begin{equation}\label{eq:W1}
 W_1(\mu,\nu)=\sup_{\operatorname{Lip}(f)\le1}
 \left|\int_\T f\dd(\mu-\nu)\right|,
\end{equation}
where the metric on $\T$ is the geodesic distance of circumference one.
For fixed irrational $M$ and degree $r$, put
\begin{equation}\label{eq:omega}
 \omega_{M,r}(\varepsilon)=
 \sup\{W_1(\mu,\lambda):\mu\text{ a probability measure on }\T,
 \ E_{M,r}(\mu)\le\varepsilon\}.
\end{equation}
The measure in this definition need not be atomic.

\begin{theorem}[Finite-frequency inversion]\label{thm:inversion}
For irrational $M>0$, every probability measure $\mu$, and every integer
$K\ge2$,
\begin{equation}\label{eq:inverse}
 W_1(\mu,\lambda)
 \le\frac8K+
 \frac{E_{M,0}(\mu)}\pi
 \sum_{n=1}^K\frac1{n|H(nM)|}
 \le\frac8K+
 \frac{E_{M,0}(\mu)(1+\log K)}{\pi h_M(K)}.
\end{equation}
In particular, $\omega_{M,r}(\varepsilon)\to0$ as $\varepsilon\downarrow0$
for every fixed irrational mesh, without a finite-type assumption.
\end{theorem}

\begin{proof}
Let $m=\lfloor K/2\rfloor$ and $L=m+1$. Normalize the square of the
Fej\'er kernel to obtain the nonnegative trigonometric polynomial
\[
 J_m(x)=\frac{F_m(x)^2}{A_m},\qquad
 A_m=\int_\T F_m^2\dd\lambda=\frac{2L^2+1}{3L}.
\]
Its degree is $2m\le K$ and its integral is one. On $|x|\le1/2$,
$F_m(x)\le L$ and $F_m(x)\le1/(4Lx^2)$. Since $A_m\ge2L/3$, splitting the
first absolute moment at $1/L$ gives
\begin{align*}
 \int_\T |x|J_m(x)\dd x
 &\le2\int_0^{1/L}x\frac{3L}{2}\dd x
   +2\int_{1/L}^{1/2}x\frac3{32L^3x^4}\dd x\\
 &\le\frac{3}{2L}+\frac3{32L}<\frac2L.
\end{align*}
Thus every one-Lipschitz function satisfies
$\norm{f-f*J_m}_\infty\le2/L$.

A periodic Lipschitz function is absolutely continuous, and integration by
parts gives $|\widehat f(n)|\le1/(2\pi|n|)$ for $n\ne0$.
Also $|\widehat J_m(n)|\le1$. With $\sigma=\mu-\lambda$, whose total
variation is at most two, we obtain
\[
 \left|\int f\dd\sigma\right|
 \le\frac4L+\frac1\pi\sum_{n=1}^K\frac{|\widehat\mu(n)|}{n}.
\]
The zeroth-moment alias identity implies
$|\widehat\mu(n)|\le E_{M,0}(\mu)/|H(nM)|$.
Finally $4/L\le8/K$ and $\sum_{n=1}^K1/n\le1+\log K$.

For the convergence assertion, first choose a fixed $K$ making $8/K$ small.
Every factor in $h_M(K)$ is nonzero. Then choose $\varepsilon$ small enough
for the second term, using $E_{M,0}\le E_{M,r}$.
\end{proof}

\begin{lemma}[Distance of a regular polygon from Haar measure]
\label[lemma]{lem:gridW}
For every $N\ge1$,
\begin{equation}\label{eq:gridW}
 W_1(\nu_N,\lambda)=\frac1{4N}.
\end{equation}
\end{lemma}
\begin{proof}
Partition the circle into the $N$ equal cells centered at the grid nodes.
Sending each cell to its center bounds every Lipschitz test difference by
$N\int_{-1/(2N)}^{1/(2N)}|x|\dd x=1/(4N)$.
For the reverse inequality use the one-Lipschitz function equal to distance
from the grid. It is zero at every atom and has Haar integral $1/(4N)$.
\end{proof}

\begin{theorem}[The inverse scale and failure of every H\"older estimate]
\label{thm:modulus}
Fix $r\ge0$.
For every irrational $M>0$, no constants $C,\alpha>0$ give a local estimate
\begin{equation}\label{eq:noholder}
 W_1(\mu,\lambda)\le C E_{M,r}(\mu)^\alpha
\end{equation}
near Haar measure. If $M$ satisfies \eqref{eq:type}, then
\begin{equation}\label{eq:omegascale}
 \omega_{M,r}(\varepsilon)
 =\exp[-\Theta(\sqrt{\log(1/\varepsilon)})].
\end{equation}
For clarity about constants, if
\[
 0<\alpha_0<\sqrt{\frac{\log2}{\tau+1}},
 \qquad \beta_0>\sqrt{2\log2},
\]
then, for all sufficiently small $\varepsilon$,
\begin{equation}\label{eq:omegaconstants}
 e^{-\beta_0\sqrt{\log(1/\varepsilon)}}
 \le\omega_{M,r}(\varepsilon)
 \le e^{-\alpha_0\sqrt{\log(1/\varepsilon)}}.
\end{equation}
\end{theorem}
\begin{proof}
For the non-H\"older assertion, use $\mu=\nu_N$.
By \cref{cor:gridupper,lem:gridW}, its error is at most
$6\,2^rG(MN)$, while its distance is $1/(4N)$.
For every fixed $\alpha>0$ the ratio of this distance to
$[6\,2^rG(MN)]^\alpha$ tends to infinity: the positive quadratic term in
$\log N$ dominates every linear term. Meanwhile $\nu_N\to\lambda$ in
$W_1$, so these counterexamples occur arbitrarily close to Haar measure.

For the finite-type upper bound write $T=\log(1/\varepsilon)$ and choose
$\alpha_1$ strictly between $\alpha_0$ and
$\sqrt{\log2/(\tau+1)}$. Put $K=\lfloor e^{\alpha_1\sqrt T}\rfloor$.
Substitute \eqref{eq:floor} into \eqref{eq:inverse}. Its second term is at
most
\[
 \exp\left[-\left(1-\frac{\tau+1}{\log2}\alpha_1^2\right)T
 +O(\sqrt T+\log T)\right],
\]
which is smaller than the first term on the stated scale. The slack between
$\alpha_0$ and $\alpha_1$ absorbs the fixed prefactors.

For the lower bound no arithmetic assumption is needed. Choose
\[
 N=\left\lceil M^{-1}
 \exp\!\left(\sqrt{2\log2\,[T+\log(6\,2^r)]}\right)\right\rceil.
\]
For large $T$, \cref{cor:gridupper} gives $E_{M,r}(\nu_N)\le\varepsilon$.
Equation \eqref{eq:gridW} then gives
$\log\omega_{M,r}(\varepsilon)\ge-\sqrt{2\log2}\sqrt T+O(1)$.
The slack in $\beta_0$ proves the lower inequality.
\end{proof}

\begin{remark}[A deterministic statement, not a statistical rate]
The error is a uniform deterministic quadrature defect, and the measure being
recovered is the phase distribution. There are no random observations or
sample-size assumptions in \cref{thm:inversion,thm:modulus}. Statistical
Wasserstein deconvolution results such as \cite{rousseau} provide context,
not a black-box proof of these particular rates.
\end{remark}

\section{Liouville meshes: arbitrarily fast approximation and slow recovery}
\label{sec:liouville}

Exact irrational rigidity alone imposes no mesh-independent lower scale for
finite positive approximation. The following theorem makes this failure
quantitative and keeps the mesh fixed after its construction.

\begin{theorem}[Arbitrary-rate obstruction]\label{thm:liouville}
Fix $r\ge0$, a nonempty open interval $I\subset(0,\infty)$, and any function
$\rho:(0,1)\to(0,\infty)$ with $\rho(\varepsilon)\to0$ as
$\varepsilon\downarrow0$. There exist a Liouville irrational $M\in I$,
integers $N_j\to\infty$, and $\varepsilon_j\downarrow0$ such that
\begin{equation}\label{eq:liouville}
 E_{M,r}(\nu_{N_j})\le\varepsilon_j,
 \qquad
 W_1(\nu_{N_j},\lambda)=\frac1{4N_j}>\rho(\varepsilon_j).
\end{equation}
Consequently $\omega_{M,r}(\varepsilon_j)>\rho(\varepsilon_j)$.
Furthermore, for any prescribed positive sequence $(\eta_N)$, the construction
can require $\varepsilon_j\le\eta_{N_j}$.
\end{theorem}
\begin{proof}
Enumerate the positive rationals as $q_1,q_2,\ldots$. Start with a closed
interval of positive length contained in $I$. Inductively choose a reduced
rational $a_j/b_j$ in the interior of the previous interval, different from
$q_1,\ldots,q_j$, with $b_j$ increasing without bound. Such rationals are
dense. Put $N_j=2^r b_j$, choosing the denominators large enough that $N_j$
is strictly increasing.

At the mesh $a_j/b_j$, the regular $N_j$-gon refines to the integer mesh
$2^r a_j$. Its integer aliases have multiplicity at least $r+1$, so its
error through degree $r$ is zero. Choose
\[
 0<\varepsilon_j<\min\{2^{-j},\varepsilon_{j-1}/2,\eta_{N_j}\}
\]
with the last condition omitted when no sequence is prescribed, and with
$\rho(\varepsilon_j)<1/(4N_j)$. The limit assumption on $\rho$ permits this.
For the first step omit $\varepsilon_{j-1}/2$.

Continuity of the fixed-grid error in the mesh gives a neighborhood of
$a_j/b_j$ throughout which that error is less than $\varepsilon_j$.
Choose a closed interval of positive length centered there, contained in
this neighborhood and the interior of the previous interval, excluding
$q_1,\ldots,q_j$, of diameter at most $2^{-j}$, and of radius less than
$b_j^{-j}$. Every one of these conditions allows a positive radius.

The nested intervals have a unique common point $M$. Exclusion of the
rationals makes $M$ irrational. Also
$0<|M-a_j/b_j|<b_j^{-j}$, so $M$ is Liouville in the usual approximation
sense. The error inequality holds for every $j$ because the final point
belongs to the $j$th interval. The distance assertion is
\cref{lem:gridW}, and the chosen inequality for $\rho$ completes the proof.
\end{proof}

The two conclusions describe the same phenomenon from opposite directions.
Exceptionally accurate rational approximations make a finite phase grid
exceptionally effective. But precisely because it remains a finite grid,
its distance from Haar measure is still $1/(4N)$; a tiny defect can therefore
carry very little information about the phase distribution.

There is no contradiction with \cref{thm:inversion}: for each constructed
irrational $M$, its inverse modulus still tends to zero. Nor does the theorem
claim that every Liouville mesh exhibits the same behavior. It constructs
meshes tailored to the prescribed obstruction.

\section{Reproducibility, finite certificates, and implementation}
\label{sec:checks}

\subsection{What can be checked exactly}

The exponent in \cref{thm:resonance} requires only integer arithmetic.
For reduced $a/b$, compute the orders in \eqref{eq:vn} for $1\le n\le N$,
take their maximum, and subtract $r$ with truncation at zero. An equivalent
constant-time arithmetic formula after integer division is
\eqref{eq:zN}; $\lfloor\log_2\lfloor N/b\rfloor\rfloor$ can be computed
from the bit length of the positive integer $\lfloor N/b\rfloor$, without
floating-point logarithms.

The supplied program compares these two routes over a finite parameter grid.
It also verifies that positivity of the exponent is exactly the inherited
support threshold \eqref{eq:oldthreshold}, and that the proposed regular
polygon never exceeds the budget. These checks are finite regression tests;
the general statements are proved in the article.

The Fej\'er coefficient sum and its squared-kernel normalization are checked
with rational arithmetic. Numerical tests of the positive and complex
Fourier witnesses supplement, but do not replace, their proofs. In
particular, no sampled collection of positive weights establishes the
universal $1/N$ inequality by itself.

\subsection{A product enclosure usable by a certified implementation}

For real $t$ and a truncation index $J$ with $|t|/2^J\le1/2$, define
\[
 P_J(t)=\prod_{h=0}^{J-1}\frac{\sin(\pi t/2^h)}{\pi t/2^h}.
\]
The tail is positive. \Cref{lem:lowerH} gives the analytic enclosure
\begin{equation}\label{eq:productcertificate}
 |P_J(t)|\exp\!\left(-\frac{4\pi^2t^2}{9\,4^J}\right)
 \le |H(t)|\le |P_J(t)|.
\end{equation}
Thus the infinite tail does not require an uncontrolled numerical cutoff.
A fully certified implementation would evaluate the finite product using
outward-rounded interval arithmetic and combine it with
\eqref{eq:productcertificate}. Derivative enclosures can instead be obtained
by finite Taylor jets and differentiated tails.

The bundled numerical code uses high-precision \texttt{mpmath}, not
directed-rounding interval arithmetic. Its evaluations are diagnostics, not
certified enclosures of the exact values. It uses the displayed analytic
tail estimate to make truncation negligible compared with the test tolerance.
The distinction between an analytic enclosure formula and a floating-point
implementation of its finite part is important.

\subsection{Recorded checks and limits}

The script records its exact assertion count, random seed, numerical precision,
parameter ranges, and residuals in \path{validation/results.json}.
It checks the Fourier upper and lower envelopes at nonintegral test points,
the leading integer-zero formula, the predicted local powers for regular
grids, and the sharp low-moment examples. A human-readable transcript is
included. The delivered run passed $358{,}664$ exact assertions and $1{,}465$
numerical diagnostics at 90 decimal digits. The largest relative residual in
the leading-zero-coefficient tests was approximately $2.29\times10^{-12}$,
using displacements of magnitude $10^{-12}$. The source contains no network
calls.

The local diagnostic tracks individual alias coefficients near rational
meshes; where a full defect is sampled, the output is explicitly labeled as
a sampled trigonometric approximation. A sampled maximum is not reported as
the true continuous supremum. No numerical result is used as a premise of a
theorem above.

\section{A formalization route and a dependency audit}
\label{sec:formal}

The article does not ship an untested Lean file labeled as a formal proof.
Instead, the following decomposition identifies concrete interfaces for a
future formalization. The repository already has the Rvachev synthesis and
alias-analysis development \cite{repoFilters,repoFixed}; its exact current
interfaces must be inspected before any reuse.

\begin{center}
\small
\begin{tabular}{p{0.25\linewidth}p{0.67\linewidth}}
\toprule
Layer & Required statement or construction\\
\midrule
Finite arithmetic & Reduced-fraction divisibility; maximum two-adic valuation
below an integer bound; the exponent formula and support threshold.\\[4pt]
Atomic moments & The annihilating polynomial identity for complex weights;
Fej\'er nonnegativity and the positive $1/N$ witness.\\[4pt]
Kernel jets & The exact zero multiplicity and leading derivative of the
infinite product; smooth parameter dependence of the alias series.\\[4pt]
Local theorem & Uniform Taylor expansion; a finite minimum of functions with
known positive leading coefficients; regular-grid refinement.\\[4pt]
Large-budget bounds & The explicit differentiated envelope and the
finite-type product floor, followed by the optimized phase-cost inequalities.\\[4pt]
Inverse theorem & The normalized squared Fej\'er kernel, its first absolute
moment bound, and the dual Wasserstein estimate.\\[4pt]
Arithmetic pathology & Nested rational intervals avoiding an enumeration,
continuity of a fixed-grid error, and the Liouville approximation property.\\
\bottomrule
\end{tabular}
\end{center}

The principal logical dependencies are short. The local optimal exponent
uses only kernel zero orders, the alias identity, a finite Fourier witness,
and Taylor's theorem. It does not depend on the finite-type hypothesis, the
Wasserstein theorem, or the Liouville construction. The positive log-square
law uses the stronger positive witness and the two envelope lemmas; it does
not depend on a claim of optimal local leading constants. The inverse theorem
uses nonzero Fourier multipliers and smoothing; its qualitative version does
not require any Diophantine inequality.

These separations are useful both mathematically and for auditing a proof
assistant implementation. In particular, no compactness of unbounded signed
measures is silently assumed. The only optimization existence statement is
for positive finite-support measures, where the simplex parameterization is
compact.

\section{Further research questions and topics}
\label{sec:future}

The questions below are proposed next projects, not assertions that every
formulation is globally open. Each needs a separate literature and repository
check before a priority claim.

\begin{question}[Optimal leading constants at a rational resonance]
For fixed $a/b,N,r$ with positive exponent $q$, does
$\Err^{\mathbb K}_{N,r}(M)/|M-a/b|^q$ have a limit as $M\to a/b$ from each
side? Determine it for $\mathbb K=+,\R,\C$.
\end{question}
\Cref{thm:resonance} gives two-sided bounds and \cref{thm:profile} gives the
constant of a regular polygon. The remaining issue is a singular optimization
over moving nodes and weights, not nonvanishing of the exponent. Small cases
such as $M_0=1$, $N=1,2,3$, and $r=0$ are concrete starting points.

\begin{question}[Uniform rational neighborhoods as the budget grows]
Find useful bounds on the constants and the radius in
\cref{thm:resonance} when $N$, $a$, or $b$ varies. Describe the crossover when
$|M-a/b|$ and $N$ tend to their limits together.
\end{question}
Nearby rational resonances can then compete. A fixed finite list of aliases
is no longer enough to obtain uniform information without additional estimates.

\begin{question}[The leading log-square constant]
Does the quotient in \eqref{eq:limconstants} converge for natural fixed meshes,
for example $M=\sqrt2$? Determine the optimal constant or its liminf and
limsup. Are regular polygons asymptotically optimal beyond the scale?
\end{question}
The present constants leave a genuine gap. The arithmetic of the entire list
$H(M),\ldots,H(NM)$ and the restrictions among a positive measure's Fourier
coefficients both matter; a minimum-multiplier bound may lose information.

\begin{question}[Signed filters at a growing budget]
Determine the asymptotic size of $\Err^{\R}_{N,r}(M)$ and
$\Err^{\C}_{N,r}(M)$ at finite-type irrational meshes. What changes under a
specified total-variation bound, allowed to depend on $N$?
\end{question}
The exact local exponent cannot improve by using signs, but that does not
settle the large-$N$ problem. The gap between $1/N$ and $1/(2^N-1)$ in the
available Fourier witnesses identifies the obstruction in the present proof.

\begin{question}[Necessary arithmetic conditions]
Characterize those irrational meshes for which the log-square complexity law,
or the stretched-logarithmic inverse scale, holds. Is a condition on a
weighted sequence of rational approximations both necessary and sufficient?
\end{question}
Finite type is sufficient and tailored Liouville examples show that some
restriction is necessary for a universal theorem. Neither statement proves
that finite type is the exact boundary.

\begin{question}[Sharper inverse constants and other metrics]
Find the optimal modulus of Haar recovery, rather than its scale. Compare
$W_p$, discrepancy, negative Sobolev norms, and recovery under smoothness or
separation assumptions on the unknown phase measure.
\end{question}
The regular-grid counterexamples rule out H\"older recovery in $W_1$ with no
extra assumptions. They do not settle every metric or every restricted model.

\begin{question}[Several dimensions and correlated filters]
For a tensor-product Rvachev kernel, determine local resonance exponents and
large-budget costs for correlated phase measures in several dimensions.
\end{question}
A zero in one coordinate can annihilate a multidimensional alias. Therefore
a direct product of the one-dimensional support thresholds need not be a
necessary bound. The finite-frequency witness should be adapted to that
geometry rather than assumed to tensor automatically.

\begin{question}[Other refinement bases and other kernels]
Which parts of the local theorem extend to smooth kernels with explicitly
known zero orders, and what replaces the log-square scale for products with
other contraction schedules?
\end{question}
The lower-bound mechanism is general: a finite moment witness converts
multiplier orders into an optimal obstruction. A matching upper construction
requires a compatible alias set, which is an additional structural condition.

\begin{question}[Certified finite-frequency implementation]
Implement interval-certified values of $b_{N,r}(M)$ and $h_M(K)$, return a
verifiable error or Wasserstein certificate, and formalize the finite
arithmetic and moment-witness layers first.
\end{question}
The formula \eqref{eq:productcertificate} supplies an explicit infinite-tail
bound. The remaining work includes robust finite-product evaluation near
multiple zeros, derivative jets, and a representation of the mesh that
supports the required arithmetic nonvanishing certificates.

\section{Conclusion}

The exact symmetry classification of phase averaging has a quantitative
extension, but not a uniformly well-conditioned one. At a fixed rational
resonance, the optimal local power is completely determined by the denominator,
the numerator's two-adic valuation, the phase budget, and the polynomial
degree. Adaptive signed or complex weights cannot improve that exponent.
An explicit regular polygon realizes it, with a leading response inherited
from a coarse half-integer mesh.

At a fixed finite-type irrational mesh, positivity yields a much stronger
large-budget statement: optimal phase errors have log-square decay, and the
corresponding phase cost is exponential in the square root of the logarithm
of the requested accuracy. The same scale governs approximate Haar recovery.
No H\"older inverse bound can hold, even locally at Haar measure. Without
arithmetic restrictions, specially constructed irrational meshes make both
approximation and inverse instability as extreme along subsequences as a
prescribed rate demands.

The remaining tasks are now more specific than the original stability
question: optimal leading constants, uniform crossover laws, signed
large-budget complexity, and the exact arithmetic boundary. The proofs and
reproducibility package separate these unresolved refinements from the results
established here.

\begin{thebibliography}{9}
\raggedright

\bibitem{arias}
J.~Arias de Reyna,
\emph{An infinitely differentiable function with compact support:
Definition and properties}, arXiv:1702.05442, 2017.
English translation of the author's article in
\emph{Revista de la Real Academia de Ciencias de Madrid} \textbf{76}
(1982), 21--38.
\url{https://arxiv.org/abs/1702.05442}.

\bibitem{rousseau}
J.~Rousseau and C.~Scricciolo,
\emph{Wasserstein convergence in Bayesian deconvolution models},
arXiv:2111.06846, 2021.
\url{https://arxiv.org/abs/2111.06846}.

\bibitem{repoRigidity}
ProveIt research archive,
\emph{Rigidity and Quantitative Limits of Phase Averaging in
Fabius--Rvachev Quadrature}, research manuscript, 28 September 2026.
Pinned repository source:
\href{https://github.com/VladimirReshetnikov/ProveIt/blob/cb1646dc442724f8e298cf259195b6c308367d80/Analysis/FabiusFunction/docs/semi-formalized-research-frontiers/drafts/inverse-and-sampling/comb-interpolation/Phase_Averaging_Rigidity/article.tex}
{\texttt{Phase\_Averaging\_Rigidity/article.tex}}.
The near-rational question occurs in its further-research section.

\bibitem{repoFilters}
ProveIt research archive,
\emph{Complete Superconvergence Phases and Optimal Phase Filters for
Rvachev Quadrature}, research manuscript, September 2026.
Pinned status and editorial notes:
\href{https://github.com/VladimirReshetnikov/ProveIt/blob/cb1646dc442724f8e298cf259195b6c308367d80/Analysis/FabiusFunction/docs/semi-formalized-research-frontiers/drafts/inverse-and-sampling/comb-interpolation/Optimal_Phase_Filters/README.md}
{\texttt{Optimal\_Phase\_Filters/README.md}}.

\bibitem{repoFixed}
ProveIt research archive,
\emph{Sharp Fixed-Phase Order of Rvachev Quadrature},
research manuscript, 29 September 2026.
Pinned repository source:
\href{https://github.com/VladimirReshetnikov/ProveIt/blob/cb1646dc442724f8e298cf259195b6c308367d80/Analysis/FabiusFunction/docs/semi-formalized-research-frontiers/drafts/inverse-and-sampling/comb-interpolation/Fixed_Phase_Order_Rvachev_Quadrature/article.tex}
{\texttt{Fixed\_Phase\_Order\_Rvachev\_Quadrature/article.tex}}.

\end{thebibliography}
\end{document}
